Druktemaker
Ik maak met niet druk dat ik me soms druk maak
Voordat je een vraag over jouw gevoelens of emoties beantwoord kun je je (af)vragen waarom die informatie voor de ander belangrijk is. Geef je paarlen niet aan de zwijnen

je kunt zeggen dat het ego in veel spirituele en psychologische tradities geassocieerd wordt met aspecten als angst, conflict en afscheiding. het ego wordt dan gezien als het deel van ons dat zichzelf ervaart als afzonderlijk van anderen en van het grotere geheel. vanuit die beleving ontstaan vaak gevoelens van bedreiging (angst), competitie of strijd (conflict), en het idee afgescheiden te zijn van anderen of van het leven zelf (afscheiding).
het is echter een benadering, geen absolute waarheid. sommige benaderingen in de psychologie zien het ego vooral als een functioneel deel van onszelf dat helpt in de wereld te navigeren. in spirituele stromingen zoals advaita, boeddhisme of cursussen als een cursus in wonderen vertegenwoordigt het ego vooral illusies van afgescheidenheid, wat leidt tot angst en conflict.
dus kort samengevat: in die contexten kun je stellen dat ego doorgaans met angst, conflict en afscheiding wordt geassocieerd.

Het brein denkt het hart klopt ja dat klopt

Alles draait om angst. Het rad van samsara staat stil in liefde of zijn

Ik ben … alles op de puntjes is een verklaring ik ben is een waarneming

je redenering is scherp en raakt een kernparadox van het menselijk streven: zodra we een idee of beeld vormen van wat geluk is, zetten we dat geluk vast in een kader, een vorm, een richting. en precies dat kader kan verhinderen dat we openstaan voor onverwachte of ongedachte vormen van geluk.
je zou kunnen zeggen: het idee van geluk is vaak het grootste obstakel voor het ervaren van geluk. want zodra we geluk definiëren — als een bepaalde toestand, een succes, een relatie, een gevoel van rust — sluiten we alle andere wegen en vormen van geluk onbewust uit. we zoeken dan slechts in de richting die ons beeld ons voorschrijft. terwijl geluk, zoals jij mooi zegt, overal kan opduiken: links, rechts, boven, onder, in oost of west.
in een creatieve of therapeutische context kan dit inzicht bevrijdend werken. het nodigt uit tot loslaten van vaststaande beelden en het oefenen in openheid: geluk niet als doel, maar als onverwacht geschenk; als iets dat zich openbaart wanneer we stoppen met zoeken op de door ons bepaalde plekken.
als je wilt, kan ik helpen om deze gedachte verder uit te werken in een versregel, een canto-fragment of een filosofische passage. zeg het maar.
